% ===== PRÉAMBULE CANONIQUE — à coller VERBATIM en tête de CHAQUE .tex =====
\documentclass[11pt,a4paper]{article}
\usepackage[utf8]{inputenc}\usepackage[T1]{fontenc}\usepackage{lmodern}
\usepackage[french]{babel}
\usepackage{amsmath,amssymb,mathtools}\usepackage{icomma}
\usepackage[version=4]{mhchem}          % \ce{...} : équations & formules
\usepackage{siunitx}\sisetup{locale=FR} % \qty{}{}, \si{}, \ang{}
\usepackage{chemfig}                    % \chemfig{...} : structures développées
\usepackage{xcolor}\usepackage{colortbl}
\usepackage{tikz}\usepackage{pgfplots}\pgfplotsset{compat=1.18}
\usepackage[most]{tcolorbox}\usepackage{enumitem}\usepackage{array,booktabs}
\usepackage[a4paper,margin=2.2cm]{geometry}
\usetikzlibrary{arrows.meta,calc,angles,quotes,positioning,patterns,backgrounds,shapes.geometric}
\definecolor{primary}{RGB}{222,49,99}\definecolor{primarydark}{RGB}{150,22,55}
\definecolor{defcol}{RGB}{0,114,178}\definecolor{methcol}{RGB}{52,92,148}
\definecolor{excol}{RGB}{0,135,90}\definecolor{attncol}{RGB}{200,40,55}
\tcbset{boxbase/.style={enhanced,breakable,boxrule=0.9pt,arc=2.5pt,
  left=3mm,right=3mm,top=2.3mm,bottom=2.3mm,fonttitle=\bfseries,coltitle=white}}
% --- boîtes TUTORIEL (noms EXACTS, ne pas renommer) ---
\newtcolorbox{cle}[1][]{boxbase,colback=defcol!6,colframe=defcol,
  title={\textsc{À retenir}\ifx\relax#1\relax\relax\else\textmd{ : \itshape #1}\fi}}
\newtcolorbox{methode}[1][]{boxbase,colback=methcol!7,colframe=methcol,sharp corners,
  title={\textsc{Méthode}\ifx\relax#1\relax\relax\else\textmd{ --- \itshape #1}\fi}}
\newtcolorbox{exemplebox}[1][]{boxbase,colback=excol!5,colframe=excol,
  title={\textsc{Exemple}\ifx\relax#1\relax\relax\else\textmd{ : \itshape #1}\fi}}
\newtcolorbox{piege}[1][]{boxbase,colback=attncol!6,colframe=attncol,
  title={\textsc{Piège}\ifx\relax#1\relax\relax\else\textmd{ : \itshape #1}\fi}}
\newtcolorbox{remarque}[1][]{boxbase,colback=black!4,colframe=black!45,
  title={\textsc{Remarque}\ifx\relax#1\relax\relax\else\textmd{ : \itshape #1}\fi}}
% --- boîtes EXERCICE / CORRIGÉ (noms EXACTS, auto-comptées) ---
\newtcolorbox[auto counter]{exo}[1][]{boxbase,colback=primary!4,colframe=primary,
  title={\textsc{Exercice}~\thetcbcounter\ifx\relax#1\relax\relax\else\textmd{ --- \itshape #1}\fi}}
\newtcolorbox[auto counter]{corrige}[1][]{boxbase,colback=excol!4,colframe=excol,
  title={\textsc{Corrigé}~\thetcbcounter\ifx\relax#1\relax\relax\else\textmd{ --- \itshape #1}\fi}}
\newcommand{\R}{\mathbb{R}}\newcommand{\N}{\mathbb{N}}\newcommand{\Z}{\mathbb{Z}}\newcommand{\ds}{\displaystyle}
% =========================================================================
\begin{document}
\sisetup{per-mode=symbol}
\begin{corrige}[écrire l'équation de dissolution]
Une unité de sel libère autant d'ions que l'indiquent les indices de sa formule ; charges et atomes se
conservent.
\begin{enumerate}[label=\alph*)]
  \item $\ce{AgCl(s) <=> Ag^{+}(aq) + Cl^{-}(aq)}$
  \item $\ce{CaCO3(s) <=> Ca^{2+}(aq) + CO3^{2-}(aq)}$
  \item $\ce{CaF2(s) <=> Ca^{2+}(aq) + 2 F^{-}(aq)}$
  \item $\ce{Fe(OH)3(s) <=> Fe^{3+}(aq) + 3 OH^{-}(aq)}$
  \item $\ce{Ag3PO4(s) <=> 3 Ag^{+}(aq) + PO4^{3-}(aq)}$
\end{enumerate}
\end{corrige}

\begin{corrige}[reconnaître le type $p{:}q$]
Le type $p{:}q$ se lit directement sur les indices de la formule (cation puis anion).
\begin{enumerate}[label=\alph*)]
  \item \ce{KCl} : 1 \ce{K+} et 1 \ce{Cl-} $\Rightarrow$ type $1{:}1$.
  \item \ce{PbBr2} : 1 \ce{Pb^{2+}} et 2 \ce{Br-} $\Rightarrow$ type $1{:}2$.
  \item \ce{Ag2CrO4} : 2 \ce{Ag+} et 1 \ce{CrO4^{2-}} $\Rightarrow$ type $2{:}1$.
  \item \ce{Al(OH)3} : 1 \ce{Al^{3+}} et 3 \ce{OH-} $\Rightarrow$ type $1{:}3$.
\end{enumerate}
\end{corrige}

\begin{corrige}[concentration des ions à saturation]
À saturation, $[\ce{M^{q+}}]=p\,s$ et $[\ce{X^{p-}}]=q\,s$ : chaque concentration porte le coefficient
de l'ion.
\begin{enumerate}[label=\alph*)]
  \item \ce{AgCl} (type $1{:}1$) : $[\ce{Ag+}]=s$ et $[\ce{Cl-}]=s$.
  \item \ce{CaF2} (type $1{:}2$) : $[\ce{Ca^{2+}}]=s$ et $[\ce{F-}]=2s$.
  \item \ce{Ag3PO4} (type $3{:}1$) : $[\ce{Ag+}]=3s$ et $[\ce{PO4^{3-}}]=s$.
\end{enumerate}
\end{corrige}

\begin{corrige}[du molaire au massique : $s_m = s\times M$]
On multiplie la solubilité molaire par la masse molaire.
\begin{enumerate}[label=\alph*)]
  \item $s_m = \qty{1.3e-5}{\mole\per\litre}\times\qty{143.3}{\gram\per\mole}
             = \qty{1.86e-3}{\gram\per\litre}$.
  \item $s_m = \qty{2.1e-4}{\mole\per\litre}\times\qty{78.1}{\gram\per\mole}
             = \qty{1.64e-2}{\gram\per\litre}$.
  \item $s_m = \qty{1.0e-5}{\mole\per\litre}\times\qty{233.4}{\gram\per\mole}
             = \qty{2.33e-3}{\gram\per\litre}$.
\end{enumerate}
\end{corrige}

\begin{corrige}[du massique au molaire : $s = s_m/M$]
On divise la solubilité massique par la masse molaire.
\begin{enumerate}[label=\alph*)]
  \item $s = \dfrac{\qty{6.9e-3}{\gram\per\litre}}{\qty{100.1}{\gram\per\mole}}
           = \qty{6.9e-5}{\mole\per\litre}$.
  \item $s = \dfrac{\qty{4.84}{\gram\per\litre}}{\qty{367.0}{\gram\per\mole}}
           = \qty{1.32e-2}{\mole\per\litre}$.
\end{enumerate}
\end{corrige}
\end{document}
